% 第9章：源域看不到目标输入的覆盖缺口（定理 thm:covariate-coverage-failure 的构造）。
% X={a,b}；源域恒产生 (a,0)，目标输入恒为 b；b 的目标标签恒为 0 或恒为 1。
% 候选 h_0(a)=h_1(a)=0，h_0(b)=0，h_1(b)=1。零一损失下目标风险：
% b 标签为 0 时 R_T(h_0)=0、R_T(h_1)=1；b 标签为 1 时 R_T(h_0)=1、R_T(h_1)=0。
\begingroup
\definecolor{FigureGreen}{HTML}{829C8B}
\definecolor{FigureGreenFill}{HTML}{E8EEE9}
\begin{tikzpicture}[figure picture,x=1mm,y=1mm]
% 左：两个输入上的概率；每行左侧为行名
\node[figure note,anchor=east] at (12,27) {源域};
\node[figure note,anchor=east] at (12,7) {目标域};
\draw[FigureStroke,line width=.6pt] (14,22) -- (48,22);
\draw[FigureStroke,line width=.6pt] (14,2) -- (48,2);
\fill[FigureBlueFill,draw=FigureBlue] (19,22) rectangle (27,34);
\node[figure note,align=center] at (23,28) {概率1\\标签0};
\draw[FigureRed,line width=.6pt,dash pattern=on 2pt off 1.5pt] (36,22) rectangle (44,34);
\node[figure note,text=FigureRed,align=center] at (40,28) {概率0};
\node[font=\figurefont\footnotesize] at (23,19.8) {$a$};
\node[font=\figurefont\footnotesize] at (40,19.8) {$b$};
\fill[FigureRedFill,draw=FigureRed] (36,2) rectangle (44,14);
\node[figure note,align=center] at (40,8) {概率1\\标签?};
\node[font=\figurefont\footnotesize] at (23,-0.2) {$a$};
\node[font=\figurefont\footnotesize] at (40,-0.2) {$b$};
% 右：两个候选在两个可能世界中的目标风险
\begin{scope}[shift={(56,0)}]
\node[figure heading,anchor=west] at (0,37) {两个候选与两个可能目标};
\foreach \x/\w/\t in {0/10/{},10/9/{$h(a)$},19/9/{$h(b)$},28/13/{$b$标0},41/13/{$b$标1}}{%
  \node[draw=FigureGrid,minimum width=\w mm,minimum height=6.4mm,inner sep=0pt,
    anchor=north west,font=\figurefont\footnotesize] at (\x,32) {\t};}
\foreach \y/\h/\pa/\pb/\ra/\rb/\fa/\fb in {25.6/$h_0$/0/0/0/1/FigureGreenFill/FigureRedFill,
                                           19.2/$h_1$/0/1/1/0/FigureRedFill/FigureGreenFill}{%
  \node[draw=FigureGrid,minimum width=10mm,minimum height=6.4mm,inner sep=0pt,anchor=north west,font=\figurefont\footnotesize] at (0,\y) {\h};
  \node[draw=FigureGrid,minimum width=9mm,minimum height=6.4mm,inner sep=0pt,anchor=north west,font=\figurefont\footnotesize] at (10,\y) {\pa};
  \node[draw=FigureGrid,minimum width=9mm,minimum height=6.4mm,inner sep=0pt,anchor=north west,font=\figurefont\footnotesize] at (19,\y) {\pb};
  \node[draw=FigureGrid,fill=\fa,minimum width=13mm,minimum height=6.4mm,inner sep=0pt,anchor=north west,font=\figurefont\footnotesize] at (28,\y) {\ra};
  \node[draw=FigureGrid,fill=\fb,minimum width=13mm,minimum height=6.4mm,inner sep=0pt,anchor=north west,font=\figurefont\footnotesize] at (41,\y) {\rb};}
\node[figure note,anchor=north west,align=left] at (0,11) {后两列为目标风险 $R_T(h)$。两种世界中\\算法收到的数据都是“全为 $(a,0)$，全为 $b$”。};
\end{scope}
\end{tikzpicture}
\endgroup
